% README.tex - Documentación detallada del proyecto CodeFlix
\documentclass[11pt,a4paper]{article}
\usepackage[utf8]{inputenc}
\usepackage[T1]{fontenc}
\usepackage[spanish]{babel}
\usepackage{hyperref}
\usepackage{geometry}
\usepackage{longtable}
\usepackage{listings}
\usepackage{tikz}
\geometry{margin=2.5cm}
\title{CodeFlix -- Documentación del Proyecto}
\author{Equipo CodeFlix}
\date{\today}

\begin{document}
\maketitle
\vspace{-1em}
\tableofcontents
\vspace{1em}

\section{Resumen ejecutivo}
CodeFlix es una aplicación web ligera que sirve como catálogo de videos/películas con una estructura de datos local basada en JSON. El repositorio incluye un servidor Node.js (archivo \texttt{server.js}), datos de ejemplo en la carpeta \texttt{data/}, assets públicos en \texttt{public/}, scripts utilitarios en \texttt{scripts/} y la interfaz estática en \texttt{src/}.

El objetivo de esta documentación es: (1) describir la arquitectura del proyecto, (2) exponer decisiones de ingeniería de software y buenas prácticas aplicadas o recomendadas, y (3) proporcionar instrucciones claras para ejecutar, mantener y ampliar el sistema.

\section{Estructura del repositorio}
\begin{longtable}{p{4cm} p{10cm}}
\textbf{Archivo/Carpeta} & \textbf{Descripción} \\\hline
\texttt{package.json} & Dependencias y scripts del proyecto (Node.js). \\
\texttt{server.js} & Servidor HTTP principal (Node/Express o http nativo según implementación). \\
\texttt{test_rating.js} & Pruebas / scripts para evaluar la puntuación o lógica (archivo de test manual). \\
\texttt{data/} & Datos de entrada (\texttt{movies.json}, \texttt{carousels.json}, \ldots). \\
\texttt{public/} & Recursos públicos servidos al cliente (imágenes, videos). \\
\texttt{scripts/} & Utilidades y scripts de normalización/transformación de datos. \\
\texttt{src/} & Código front-end (HTML, CSS, JS). \\
\texttt{README.tex} & Esta documentación. \\
\end{longtable}

\section{Componentes principales}
\subsection{Servidor}
El servidor gestiona rutas estáticas y API mínimas (por ejemplo, servir \texttt{/data/movies.json} o endpoints personalizados). Está pensado para desarrollo local y prototipado.

\subsection{Datos}
Los datos están en archivos JSON en \texttt{data/}. Estos contienen los catálogos y metadatos de las películas. Existen scripts en \texttt{scripts/} para normalizar o enriquecer dichos JSON.

\subsection{Frontend}
El código en \texttt{src/} provee las vistas: una página index (catálogo) y una página de detalle (\texttt{movie.html}). Los scripts ubicados en \texttt{src/scripts/} implementan el comportamiento en cliente.

\section{Arquitectura y flujo de datos}
\subsection{Visión por capas}
CodeFlix sigue una arquitectura en capas simple y clásica:
- Capa de datos: archivos JSON en \texttt{data/}. 
- Capa de servidor: \texttt{server.js} expone datos y sirve los activos. 
- Capa de presentación: HTML/CSS/JS en \texttt{src/}. 

El flujo básico es: cliente solicita páginas estáticas -> servidor devuelve HTML/CSS/JS -> el JS del cliente solicita los JSON (p. ej. \texttt{/data/movies.json}) -> el cliente renderiza la interfaz.

\subsection{Diagrama (simplificado)}
\begin{center}
\begin{tikzpicture}[node distance=2.2cm, auto]
  \node (browser) [draw, rectangle] {Navegador (UI)};
  \node (server) [draw, rectangle, right of=browser, node distance=4.5cm] {Servidor Node.js (\texttt{server.js})};
  \node (data) [draw, rectangle, below of=server] {Datos: \texttt{data/*.json}};
  \node (public) [draw, rectangle, above of=server] {Assets: \texttt{public/}};
  \draw[->] (browser) to node {HTTP requests} (server);
  \draw[->] (server) to node {serve static} (public);
  \draw[->] (server) to node {read/serve JSON} (data);
\end{tikzpicture}
\end{center}

\section{Decisiones de ingeniería de software}
\subsection{Modularidad y separación de responsabilidades}
El proyecto está dividido en módulos por responsabilidad: servidor, datos, utilidades y presentación. Esta separación facilita pruebas unitarias, sustitución de la capa de persistencia y despliegue independiente.

\subsection{Patrones sugeridos}
- API REST simple para exponer recursos (movies, carousels, contacts). 
- Modelo Controlador (handler) para la lógica de endpoints en \texttt{server.js}. 
- Scripts de normalización en \texttt{scripts/} para ETL y mantenimiento de datos.

\subsection{Mantenibilidad}
- Mantener contratos estables en los JSON (campos, tipos) reduce roturas en cliente.
- Documentar las estructuras de los objetos de datos (p. ej. esquema para una entrada de \texttt{movies.json}).

\subsection{Pruebas y calidad}
- Integrar un runner de pruebas (Jest / Mocha) para verificar lógica de transformación y rutas.
- Añadir linting (ESLint) y formateo automático (Prettier) como paso de CI.

\subsection{Escalabilidad y producción}
- Para producción, reemplazar archivos JSON por una base de datos (SQLite, PostgreSQL o MongoDB) si el volumen crece.
- Usar un servidor estático (nginx) frente al servidor node para entregar assets y balanceo.

\subsection{Seguridad}
- Validar y sanitizar cualquier entrada (si se añaden endpoints que reciban datos).
- Evitar exponer paths de archivos locales si se añaden upload/download.

\section{Instrucciones de desarrollo y ejecución}
\subsection{Requisitos}
- Node.js v14+ (recomendado). 

\subsection{Instalación}
Abrir terminal en la raíz del proyecto y ejecutar:
\begin{lstlisting}
npm install
\end{lstlisting}

\subsection{Ejecución en entorno de desarrollo}
\begin{lstlisting}
node server.js
\end{lstlisting}
Luego abrir en el navegador \texttt{http://localhost:3000} (o el puerto configurado en \texttt{server.js}).

\section{Documentación del código y tareas realizadas}
En el repositorio se han añadido comentarios y JSDoc en los puntos de entrada principales (servidor, scripts de normalización y scripts cliente) para facilitar la lectura y el mantenimiento. Mantener comentarios concisos, actualizados y centrados en la intención y contrato de las funciones.

\section{Cómo contribuir}
- Crea una rama con prefijo \texttt{feature/} o \texttt{fix/}. 
- Añade tests para cambios lógicos. 
- Abre un pull request descriptivo y referencia el issue si existe.

\section{Apéndice: Esquema sugerido para \texttt{movies.json}}
Ejemplo de campos básicos que deben documentarse y mantenerse estables:
\begin{itemize}
  \item \texttt{id}: identificador único (string o número)
  \item \texttt{title}: título de la película (string)
  \item \texttt{year}: año de lanzamiento (número)
  \item \texttt{duration}: duración en minutos (número opcional)
  \item \texttt{tags}: lista de etiquetas/categorías (array de strings)
  \item \texttt{videoUrl}: ruta o URL del recurso de video (string)
  \item \texttt{poster}: ruta a la imagen/póster (string)
\end{itemize}

\section{Contacto}
Para preguntas o coordinación, abrir un issue en el repositorio o contactar al autor del proyecto.

\end{document}
